%! Sets of Vectors and the Road to Subspaces — Unit 3, Lesson 3.0 — Guided Notes (Answer Key)
\documentclass[10pt]{article}
\usepackage{linalg-article}
\usepackage{linalg-key}   % pulls in -boxes; do NOT also load -boxes
\newcommand{\TallMath}[1]{$\displaystyle #1\rule[-1.4em]{0pt}{3.2em}$}
\newcommand{\vv}{\mathbf{v}}
\newcommand{\ww}{\mathbf{w}}
\newcommand{\zero}{\mathbf{0}}

\begin{document}

\pageheader{Unit 3, Lesson 3.0}{Guided Notes --- Answer Key}
\namedateperiod

\vspace{0.12in}

\begin{objectivebox}
\textbf{By the end of this lesson, I will be able to\dots}
\begin{itemize}\setlength{\itemsep}{2pt}
  \item describe the \emph{span} of some vectors as a line or plane \emph{through the origin};
  \item test whether a set of vectors is a \emph{subspace} --- closed under $+$ and scaling;
  \item recognize the \emph{column space} $C(A)$ as a subspace and see where Unit~3 is headed.
\end{itemize}
\end{objectivebox}

\vspace{0.10in}

\begin{vocabbox}
\small
Fill in each term as we build it together.
\vspace{2pt}
\vocabans{Vector space $\mathbb{R}^n$}{the set of all vectors with $n$ real components ($\mathbb{R}^2$ is the plane, $\mathbb{R}^3$ is all of space)}
\vocabans{Span (of some vectors)}{the set of all their linear combinations $c\vv+d\ww$ --- a line, plane, or all of $\mathbb{R}^n$, through the origin}
\vocabans{Subspace}{a set of vectors closed under addition and scaling (so it contains $\zero$)}
\vocabans{Closed under addition / scaling}{adding two members, or scaling a member by any $c$, always stays inside the set}
\vocabans{Column space $C(A)$}{the span of the columns of $A$ = all reachable targets $A\mathbf x$; a subspace}
\end{vocabbox}

\begin{hookbox}
A drone starts at the origin and repeats two moves, forwards or backwards as often as it likes:
$\vv=\begin{bsmallmatrix} 2\\1 \end{bsmallmatrix}$ and $\ww=\begin{bsmallmatrix} 1\\3 \end{bsmallmatrix}$.
\begin{itemize}\setlength{\itemsep}{1pt}
  \item Which points can it reach with $c\vv+d\ww$? Describe the set. \ansline{The whole plane $\mathbb{R}^2$ --- the two moves point in different directions.}
  \item Now $\ww$ becomes $\begin{bsmallmatrix} 4\\2 \end{bsmallmatrix}=2\vv$. What can it reach now? \ansline{Only the line through the origin in the direction of $\vv$.}
\end{itemize}
\end{hookbox}

\begin{notesbox}{1. The Arena: $\mathbb{R}^n$ and Span}
The \textbf{vector space $\mathbb{R}^n$} is the set of \emph{all} vectors with $n$ real
components. So $\mathbb{R}^2$ is the whole \ans{plane}, and $\mathbb{R}^3$ is all of space.

\vspace{3pt}
From Lessons 1.1 and 1.3, the \textbf{span} of some vectors is the set of \emph{all} their linear
combinations $c\vv + d\ww$. Geometrically the span is one of:
\begin{itemize}[leftmargin=1.5em, itemsep=1pt]
  \item a \textbf{line} through the origin --- when the vectors point along \ans{the same direction};
  \item a \textbf{plane} through the origin --- when they point in \ans{different (independent)} directions;
  \item \textbf{all of $\mathbb{R}^n$}.
\end{itemize}
Every span passes through the origin, because taking all weights $=0$ gives the \ans{zero}
vector.
\end{notesbox}

\begin{notesbox}{2. What Makes a Set a Subspace?}
A \textbf{subspace} is a set $S$ of vectors that you cannot \emph{escape} by combining its
members. Two rules --- the \textbf{closure test}:
\begin{enumerate}[leftmargin=1.9em, itemsep=3pt]
  \item[\textbf{(a)}] \textbf{Closed under addition:} if $\vv$ and $\ww$ are in $S$, then
        $\vv+\ww$ is \ans{in} $S$.
  \item[\textbf{(b)}] \textbf{Closed under scaling:} if $\vv$ is in $S$, then $c\vv$ is in $S$ for
        \emph{every} number $c$.
\end{enumerate}
Rule (b) with $c=0$ forces the \textbf{zero vector} $\zero$ to be in $S$. So a quick first check:
\emph{if a set does not contain $\zero$, it is} \ans{never} \emph{a subspace}.
\end{notesbox}

\begin{notesbox}{3. A Catalog You Can See}
In $\mathbb{R}^2$, compare three sets (each figure is the shaded/heavy set):
\begin{center}
\begin{tikzpicture}[scale=0.42]
  % Panel 1: line through origin  (subspace)
  \begin{scope}
    \draw[step=1, linegray!45, thin] (-3.2,-3.2) grid (3.2,3.2);
    \draw[->, thick, charcoal] (-3.4,0)--(3.4,0);
    \draw[->, thick, charcoal] (0,-3.4)--(0,3.4);
    \draw[very thick, burgundy] (-3.0,-1.5)--(3.0,1.5);
    \fill[black] (0,0) circle (3pt);
    \node[charcoal, font=\scriptsize] at (0,-4.1){line through $\zero$};
    \node[burgundy, font=\scriptsize\bfseries] at (0,-4.9){SUBSPACE};
  \end{scope}
  % Panel 2: line NOT through origin  (not a subspace)
  \begin{scope}[xshift=9.2cm]
    \draw[step=1, linegray!45, thin] (-3.2,-3.2) grid (3.2,3.2);
    \draw[->, thick, charcoal] (-3.4,0)--(3.4,0);
    \draw[->, thick, charcoal] (0,-3.4)--(0,3.4);
    \draw[very thick, royalblue] (-3.0,-0.5)--(2.5,2.25);
    \fill[black] (0,0) circle (3pt);
    \node[charcoal, font=\scriptsize] at (0,-4.1){line missing $\zero$};
    \node[royalblue, font=\scriptsize\bfseries] at (0,-4.9){NOT a subspace};
  \end{scope}
  % Panel 3: first quadrant  (not a subspace)
  \begin{scope}[xshift=18.4cm]
    \fill[burgundy!18] (0,0) rectangle (3.2,3.2);
    \draw[step=1, linegray!45, thin] (-3.2,-3.2) grid (3.2,3.2);
    \draw[->, thick, charcoal] (-3.4,0)--(3.4,0);
    \draw[->, thick, charcoal] (0,-3.4)--(0,3.4);
    \fill[black] (0,0) circle (3pt);
    \node[charcoal, font=\scriptsize] at (0,-4.1){first quadrant};
    \node[royalblue, font=\scriptsize\bfseries] at (0,-4.9){NOT a subspace};
  \end{scope}
\end{tikzpicture}
\end{center}

\vspace{2pt}
Fill in \emph{why} each middle/right set fails the closure test:
\begin{itemize}[leftmargin=1.5em, itemsep=2pt]
  \item \textbf{Line missing $\zero$:} it does not contain \ans{$\zero$}; and scaling a point on
        it by $0$ lands at the origin, which is \emph{off} the line. \textbf{Fails.}
  \item \textbf{First quadrant:} it \emph{is} closed under adding, but scaling
        $\begin{bsmallmatrix} 1\\1 \end{bsmallmatrix}$ by $c=\ans{$-1$}$ gives
        $\begin{bsmallmatrix} -1\\-1 \end{bsmallmatrix}$, which \emph{leaves} the quadrant.
        \textbf{Fails.}
\end{itemize}
\textbf{The complete list in $\mathbb{R}^2$:} just $\{\zero\}$, any line through the origin, and
all of $\mathbb{R}^2$. (In $\mathbb{R}^3$, add every \ans{plane} through the origin.) Every
subspace is a flat piece \emph{pinned to the origin}.
\end{notesbox}

\begin{notesbox}{4. The Bridge: The Column Space Is a Subspace}
The \textbf{column space} $C(A)$ is the span of the columns of $A$ = every target $A\mathbf x$ you
can \textbf{reach}. Why is it a subspace? Add two reachable targets, or scale one, and the result
is still a combination of the columns --- still \ans{reachable}. So $C(A)$ passes the closure test.

\vspace{3pt}
This renames a Unit~2 idea: ``$A\mathbf x=\mathbf b$ has a solution exactly when $\mathbf b$ is
reachable'' becomes ``\dots exactly when $\mathbf b$ is in the subspace \ans{$C(A)$}.''

\vspace{3pt}
\textbf{The road ahead (Unit 3).} \textbf{3.1} vector spaces \& subspaces\; $\bullet$\;
\textbf{3.2} the \emph{nullspace} $N(A)$: all solutions of $A\mathbf x=\zero$ (another subspace)\;
$\bullet$\; \textbf{3.3} the complete solution to $A\mathbf x=\mathbf b$\; $\bullet$\; \textbf{3.4}
independence, basis \& dimension\; $\bullet$\; \textbf{3.5} the dimensions of the four subspaces.
\end{notesbox}

\begin{practicebox}
Decide whether each set is a subspace of $\mathbb{R}^2$. If not, give a specific combination that
\emph{escapes} it.
\begin{enumerate}[leftmargin=1.6em, itemsep=4pt]
  \item All vectors on the line $y=3x$ (i.e.\ all multiples of $\begin{bsmallmatrix} 1\\3 \end{bsmallmatrix}$). \ansline{Subspace --- a line through the origin; sums and scalings of multiples of $\begin{bsmallmatrix} 1\\3 \end{bsmallmatrix}$ stay multiples of it.}
  \item All vectors on the line $y=3x+2$. \ansline{Not a subspace --- it misses $\zero$ (at $x=0$, $y=2$). E.g.\ $2\cdot\begin{bsmallmatrix} 0\\2 \end{bsmallmatrix}=\begin{bsmallmatrix} 0\\4 \end{bsmallmatrix}$ is off the line.}
  \item All of $\mathbb{R}^2$. \ansline{Subspace --- every sum and scaling of $2$-vectors is again a $2$-vector.}
\end{enumerate}
\end{practicebox}

\begin{teachernote}
The single new idea is \emph{closure}, and the single most common error is forgetting the origin.
Keep returning to ``is $\zero$ in it?'' as the fast reject, then ``can I scale out of it?'' as the
sharper test (the first-quadrant trap). Section~4 is the payoff: $C(A)$ is the first named
subspace, and 3.2's nullspace is the second --- flag both so Unit~3 feels like one story.
\end{teachernote}

\end{document}
